\begin{figure*}[t]
\centering

\resizebox{\textwidth}{!}{%
\begin{tikzpicture}[
    font=\sffamily,
    >=Latex,
    dataBox/.style={
        draw=wfBlue,
        fill=wfBlueLight,
        line width=0.8pt,
        rounded corners=2pt,
        text width=4.5cm,
        minimum height=3.1cm,
        inner sep=7pt,
        align=left
    },
    patentBox/.style={
        draw=wfOrange,
        fill=wfOrangeLight,
        line width=0.8pt,
        rounded corners=2pt,
        text width=5.0cm,
        minimum height=3.1cm,
        inner sep=7pt,
        align=left
    },
    policyBox/.style={
        draw=wfPurple,
        fill=wfPurpleLight,
        line width=0.8pt,
        rounded corners=2pt,
        text width=5.0cm,
        minimum height=3.1cm,
        inner sep=7pt,
        align=left
    },
    analysisBox/.style={
        draw=wfTeal,
        fill=wfTealLight,
        line width=0.8pt,
        rounded corners=2pt,
        text width=5.0cm,
        minimum height=3.0cm,
        inner sep=7pt,
        align=left
    },
    compareBox/.style={
        draw=wfTealDark,
        fill=white,
        line width=0.9pt,
        rounded corners=2pt,
        text width=5.0cm,
        minimum height=3.0cm,
        inner sep=7pt,
        align=left
    },
    outputBox/.style={
        draw=wfOrangeDark,
        fill=wfOrangeLight,
        line width=1pt,
        rounded corners=2pt,
        text width=7.2cm,
        minimum height=2.8cm,
        inner sep=7pt,
        align=center
    },
    blueArrow/.style={
        ->,
        draw=wfBlueDark,
        line width=1.3pt
    },
    patentArrow/.style={
        ->,
        draw=wfOrangeDark,
        line width=1.3pt
    },
    policyArrow/.style={
        ->,
        draw=wfPurpleDark,
        line width=1.3pt
    },
    analysisArrow/.style={
        ->,
        draw=wfTealDark,
        line width=1.3pt
    },
    outputArrow/.style={
        ->,
        draw=wfOrangeDark,
        line width=1.3pt
    }
]

% ============================================================
% BLOCK A — SCHOLARLY CORPUS CONSTRUCTION
% ============================================================

\node[
    anchor=west,
    fill=wfBlue,
    text=white,
    rounded corners=2pt,
    font=\bfseries\small,
    inner xsep=9pt,
    inner ysep=5pt
] (blockA) at (-2.5,2.1)
{BLOCK A --- SCHOLARLY CORPUS CONSTRUCTION};


\node[dataBox] (scopus) at (0,0)
{
    {\color{wfBlueDark}\bfseries\large A1}
    \hfill
    {\color{wfGray}\scriptsize SEARCH}

    \vspace{3pt}

    {\bfseries Scopus Literature Retrieval}

    \vspace{6pt}

    {\scriptsize
    \textbf{Database:} Scopus

    \vspace{3pt}

    ML--OPF scholarly publications are identified using the
    predefined literature-search strategy.

    \vspace{5pt}

    Titles, abstracts, publication years, authors,
    keywords, and DOIs are retained.
    }
};


\node[dataBox, right=0.65cm of scopus] (clean)
{
    {\color{wfBlueDark}\bfseries\large A2}
    \hfill
    {\color{wfGray}\scriptsize CLEAN}

    \vspace{3pt}

    {\bfseries Data Cleaning}

    \vspace{6pt}

    {\scriptsize
    Records are checked for:

    \vspace{4pt}

    $\bullet$ duplicate publications

    \vspace{2pt}

    $\bullet$ title and abstract availability

    \vspace{2pt}

    $\bullet$ publication year

    \vspace{2pt}

    $\bullet$ DOI availability
    }
};


\node[dataBox, right=0.65cm of clean] (doi)
{
    {\color{wfBlueDark}\bfseries\large A3}
    \hfill
    {\color{wfGray}\scriptsize DOI}

    \vspace{3pt}

    {\bfseries DOI Identification}

    \vspace{6pt}

    {\scriptsize
    DOIs provide the common identifier used to link scholarly
    publications across Scopus, Lens, and Overton.

    \vspace{7pt}

    {\color{wfBlueDark}\bfseries
    Common ML--OPF scholarly corpus}
    }
};


\draw[blueArrow]
(scopus.east) -- (clean.west);

\draw[blueArrow]
(clean.east) -- (doi.west);


% ============================================================
% BLOCK B — PATENT AND POLICY CITATION PATHWAYS
% ============================================================

\node[
    anchor=west,
    fill=wfGray,
    text=white,
    rounded corners=2pt,
    font=\bfseries\small,
    inner xsep=9pt,
    inner ysep=5pt
] (blockB) at (-2.5,-4.1)
{BLOCK B --- PATENT AND POLICY CITATION PATHWAYS};


\node[
    patentBox,
    below left=2.3cm and 0.20cm of doi
] (lens)
{
    {\color{wfOrangeDark}\bfseries\large B1}
    \hfill
    {\color{wfGray}\scriptsize PATENTS}

    \vspace{3pt}

    {\bfseries Lens Patent-Citation Pathway}

    \vspace{6pt}

    {\scriptsize
    \textbf{Database:} Lens

    \vspace{3pt}

    Scopus DOIs identify scholarly publications cited by patents.

    \vspace{7pt}

    {\color{wfOrangeDark}\bfseries
    Patent-cited scholarly publications}

    \vspace{2pt}

    {\color{wfOrangeDark}\bfseries
    $N=778$}
    }
};


\node[
    policyBox,
    below right=2.3cm and 0.20cm of doi
] (overton)
{
    {\color{wfPurpleDark}\bfseries\large B2}
    \hfill
    {\color{wfGray}\scriptsize POLICY}

    \vspace{3pt}

    {\bfseries Overton Policy-Citation Pathway}

    \vspace{6pt}

    {\scriptsize
    \textbf{Database:} Overton

    \vspace{3pt}

    Scopus DOIs identify scholarly publications cited by policy documents.

    \vspace{7pt}

    {\color{wfPurpleDark}\bfseries
    Policy-cited scholarly publications}

    \vspace{2pt}

    {\color{wfPurpleDark}\bfseries
    $N=392$}
    }
};


\draw[
    patentArrow,
    rounded corners=5pt
]
(doi.south)
-- ++(0,-0.65cm)
-| (lens.north);


\draw[
    policyArrow,
    rounded corners=5pt
]
(doi.south)
-- ++(0,-0.65cm)
-| (overton.north);


% ============================================================
% BLOCK C — CORPUS-SPECIFIC TOPIC MODELING
% ============================================================

\node[
    anchor=west,
    fill=wfTeal,
    text=white,
    rounded corners=2pt,
    font=\bfseries\small,
    inner xsep=9pt,
    inner ysep=5pt
] (blockC) at (-2.5,-10.2)
{BLOCK C --- CORPUS-SPECIFIC TOPIC MODELING};


\node[
    analysisBox,
    below=2.1cm of lens
] (patentdtm)
{
    {\color{wfTealDark}\bfseries\large C1}
    \hfill
    {\color{wfGray}\scriptsize DTM}

    \vspace{3pt}

    {\bfseries Patent-Cited Corpus}

    \vspace{6pt}

    {\scriptsize
    Identical abstract preprocessing.

    \vspace{4pt}

    \textbf{Final DTM}

    778 documents

    1,295 terms

    Minimum DF = 8
    }
};


\node[
    analysisBox,
    below=2.1cm of overton
] (policydtm)
{
    {\color{wfTealDark}\bfseries\large C2}
    \hfill
    {\color{wfGray}\scriptsize DTM}

    \vspace{3pt}

    {\bfseries Policy-Cited Corpus}

    \vspace{6pt}

    {\scriptsize
    Identical abstract preprocessing.

    \vspace{4pt}

    \textbf{Final DTM}

    392 documents

    1,371 terms

    Minimum DF = 4
    }
};


\draw[analysisArrow]
(lens.south) -- (patentdtm.north);

\draw[analysisArrow]
(overton.south) -- (policydtm.north);


\node[
    analysisBox,
    below=0.75cm of patentdtm
] (patentlda)
{
    {\color{wfTealDark}\bfseries\large C3}
    \hfill
    {\color{wfGray}\scriptsize LDA}

    \vspace{3pt}

    {\bfseries Patent-Cited LDA}

    \vspace{6pt}

    {\scriptsize
    Independent topic model.

    \vspace{4pt}

    Topic-number selection:

    $\bullet$ held-out perplexity

    $\bullet$ semantic coherence

    $\bullet$ topic similarity

    $\bullet$ interpretability
    }
};


\node[
    analysisBox,
    below=0.75cm of policydtm
] (policylda)
{
    {\color{wfTealDark}\bfseries\large C4}
    \hfill
    {\color{wfGray}\scriptsize LDA}

    \vspace{3pt}

    {\bfseries Policy-Cited LDA}

    \vspace{6pt}

    {\scriptsize
    Independent topic model.

    \vspace{4pt}

    Topic-number selection:

    $\bullet$ held-out perplexity

    $\bullet$ semantic coherence

    $\bullet$ topic similarity

    $\bullet$ interpretability
    }
};


\draw[analysisArrow]
(patentdtm.south) -- (patentlda.north);

\draw[analysisArrow]
(policydtm.south) -- (policylda.north);


\node[
    analysisBox,
    below=0.75cm of patentlda,
    minimum height=2.0cm
] (patenttopics)
{
    {\color{wfTealDark}\bfseries\large C5}
    \hfill
    {\color{wfGray}\scriptsize TOPICS}

    \vspace{3pt}

    {\bfseries Patent Topics}

    \vspace{4pt}

    {\scriptsize
    Corpus-specific latent themes
    }
};


\node[
    analysisBox,
    below=0.75cm of policylda,
    minimum height=2.0cm
] (policytopics)
{
    {\color{wfTealDark}\bfseries\large C6}
    \hfill
    {\color{wfGray}\scriptsize TOPICS}

    \vspace{3pt}

    {\bfseries Policy Topics}

    \vspace{4pt}

    {\scriptsize
    Corpus-specific latent themes
    }
};


\draw[analysisArrow]
(patentlda.south) -- (patenttopics.north);

\draw[analysisArrow]
(policylda.south) -- (policytopics.north);


% ============================================================
% BLOCK D — CROSS-CORPUS INTERPRETATION
% ============================================================

\node[
    anchor=west,
    fill=wfTealDark,
    text=white,
    rounded corners=2pt,
    font=\bfseries\small,
    inner xsep=9pt,
    inner ysep=5pt
] (blockD) at (-2.5,-23.2)
{BLOCK D --- CROSS-CORPUS INTERPRETATION \& COMPARISON};


\node[
    compareBox,
    below=2.0cm of patenttopics
] (interpret)
{
    {\color{wfTealDark}\bfseries\large D1}
    \hfill
    {\color{wfGray}\scriptsize INTERPRET}

    \vspace{3pt}

    {\bfseries Manual Interpretation}

    \vspace{6pt}

    {\scriptsize
    $\bullet$ highest-probability terms

    \vspace{2pt}

    $\bullet$ representative publications

    \vspace{2pt}

    $\bullet$ domain-informed validation
    }
};


\node[
    compareBox,
    right=0.65cm of interpret
] (mapping)
{
    {\color{wfTealDark}\bfseries\large D2}
    \hfill
    {\color{wfGray}\scriptsize MAPPING}

    \vspace{3pt}

    {\bfseries Cross-Corpus Mapping}

    \vspace{6pt}

    {\scriptsize
    Related patent and policy topics are organized into a
    common thematic framework.

    \vspace{5pt}

    Topic numbers are not assumed to correspond directly.
    }
};


\node[
    compareBox,
    right=0.65cm of mapping
] (prevalence)
{
    {\color{wfTealDark}\bfseries\large D3}
    \hfill
    {\color{wfGray}\scriptsize PREVALENCE}

    \vspace{3pt}

    {\bfseries Thematic Prevalence}

    \vspace{6pt}

    {\scriptsize
    Document--topic probabilities are aggregated across
    mapped themes.

    \vspace{5pt}

    Patent--policy differences are expressed in
    percentage points.
    }
};


\node[
    compareBox,
    right=0.65cm of prevalence
] (temporal)
{
    {\color{wfTealDark}\bfseries\large D4}
    \hfill
    {\color{wfGray}\scriptsize TEMPORAL}

    \vspace{3pt}

    {\bfseries Temporal Comparison}

    \vspace{6pt}

    {\scriptsize
    Annual probability-weighted prevalence is examined
    together with annual publication coverage.

    \vspace{5pt}

    Theme evolution is compared over time.
    }
};


\draw[
    analysisArrow,
    rounded corners=5pt
]
(patenttopics.south)
-- ++(0,-0.65cm)
-| (interpret.north);


\draw[
    analysisArrow,
    rounded corners=5pt
]
(policytopics.south)
-- ++(0,-0.65cm)
-| (interpret.north);


\draw[analysisArrow]
(interpret.east) -- (mapping.west);

\draw[analysisArrow]
(mapping.east) -- (prevalence.west);

\draw[analysisArrow]
(prevalence.east) -- (temporal.west);


% ============================================================
% FINAL OUTPUT
% ============================================================

\node[
    outputBox,
    below=0.85cm of mapping
] (output)
{
    {\color{wfOrangeDark}\bfseries
    COMPARATIVE KNOWLEDGE-TRANSLATION PROFILE}

    \vspace{7pt}

    {\scriptsize
    \textbf{Technology-facing themes}
    \qquad
    $\bullet$
    \qquad
    \textbf{Shared / bridge themes}

    \vspace{4pt}

    \textbf{Policy-facing themes}
    }
};


\draw[
    outputArrow,
    rounded corners=4pt
]
(temporal.south)
-- ++(0,-0.45cm)
-| (output.east);


\end{tikzpicture}%
}

\caption{
Research workflow for the comparative analysis of patent-cited and
policy-cited ML--OPF scholarly literature.
Block A constructs the underlying Scopus scholarly corpus and identifies
DOIs for cross-database linkage.
Block B identifies patent-cited and policy-cited scholarly publications
through Lens and Overton.
Block C applies identical text preprocessing and corpus-specific
document-term matrix construction before estimating independent LDA
models.
Block D interprets the resulting topics, maps related topics into a
common thematic framework, and compares their probability-weighted
prevalence and temporal evolution.
}

\label{fig:research_workflow}

\end{figure*}
